\documentclass[11pt]{article}
\usepackage{amsmath}
\usepackage{booktabs}
\title{Paper A Skeleton: Identified Component-Box Diagnostics for Low-Ell Anisotropy}
\author{HTT Collaboration Draft}
\date{Generated diagnostic skeleton; not a compiled-submission claim}
\begin{document}
\maketitle

\begin{abstract}
We present a claim-tiered identifiability and diagnostic framework. The current
artifacts are diagnostic and transfer-conditional; they do not use native low-ell
solver output and do not promote morphology-family conclusions.
\end{abstract}

\section{Component-Box Identified Sets}
The component vector is bounded by signed lower and upper boxes. The curvature
component is reported with both open and all-branch intervals.

\section{Two-Stage Thresholds}
Known-covariance chi-square thresholds and simulation-estimated covariance
Hotelling/F thresholds are separate policies with explicit metadata.

\section{Depth-Gap Interval Propagation}
The joint interval is a subset of the naive quotient interval. Strictness requires
a shared-extrema conflict; aligned shared boxes can collapse to equality.

\section{Diagnostic Data Appendix}
The K5/CF4 card is a pipeline-closure diagnostic while Sigma2, W2, and Omega-k
components remain blocked placeholders. Report figures are appendix diagnostics
unless matched-null and covariance gates pass.

\end{document}
